% =====================================================================================
% Report outline v2 (2026-11-27), updated from the original Overleaf draft to match
% the results in project_star_notes.txt sections 9-14 and figures/ (make_figures.py).
%
% DEADLINES (course page): slides due Mon 1 Dec 11:59 AM; 8-page report due Mon 1 Dec
% 11:59 PM; 8-minute presentation Tue 2 Dec in class.
%
% Page budget for 8 pages (the exemplar puts References + Contributions after the
% 8 pages; check the course page if unsure):
%   1 Introduction ............................. 0.75
%   2 Project STAR Data (with Table 1) ......... 0.75
%   3 Methods .................................. 1.75
%   4 Experiments and Results (3 figs, 1 table)  3.5
%   5 Discussion ............................... 0.9
%   6 Conclusion ............................... 0.25
% Main-text figures, in priority order: Fig 2 (frontier), Fig 4 (gender shares),
% Fig 5 (trees), Fig 3 (capacity). Fig 1 (gain by group) only if space allows.
%
% Every number below is from notes sections 9-14 or figures/*.csv. Notes in
% "% NOTE:" comments are for us, not for the report.
%
% FIGURES: upload the PNGs from figures/ into an Overleaf folder called figures/.
% Until then each figure shows a placeholder box.
% =====================================================================================

\documentclass[11pt]{article}
\usepackage[letterpaper,margin=0.95in]{geometry}
\usepackage{amsmath,amssymb,booktabs,graphicx}
\usepackage[hidelinks]{hyperref}
\usepackage{enumitem}
\setlist{nosep}
\newcommand{\placeholder}[1]{\par\noindent\textit{[To write: #1]}\par}
\newcommand{\treatmentset}{\mathcal{T}}
\newcommand{\fig}[3]{% file, width, fallback text
  \IfFileExists{#1}{\includegraphics[width=#2]{#1}}%
  {\fbox{\parbox[c][1.3in][c]{0.8\linewidth}{\centering Upload #1 (#3)}}}}

\title{Optimizing Classroom Resource Allocation Using Optimal Policy Trees}
\author{Adi Palaniappan \and Daniel Krupp}
\date{December 2026}

\begin{document}
\maketitle

% =====================================================================================
\section{Introduction}  % ~0.75 page

\subsection{Motivation and Background}
\placeholder{Condense the proposal's motivation (proposal.pdf p.1). Project STAR:
randomized Tennessee experiment, about 11,600 students in 79 schools, K--3; three class
types (regular 22--25 students, small 13--17, regular with a full-time aide). Literature:
small classes about +0.22 SD after four years \cite{krueger1999}, roughly three months of
schooling \cite{whitehurst2011}; little average effect of aides. Small classes and aides
cost staff and rooms, so schools can only offer a limited number of places. STAR tells us
whether small classes help on average, not who should get the limited places.}
% NOTE: Krueger reported larger gains for free-lunch and Black students. In our
% kindergarten data the gain is the same with and without free lunch (0.215 vs 0.212).
% Keep literature claims and our findings separate.

\subsection{Research Question and Contributions}
\placeholder{Question: can a shallow, interpretable rule using student characteristics
decide who gets limited small-class and aide places, and raise expected achievement
compared with simple rules, under school-level capacity limits?}
\placeholder{Contributions (short bullet list):
(i) a two-stage pipeline: cross-fitted doubly robust rewards, then Optimal Policy Trees,
evaluated on held-out schools;
(ii) per-school capacity limits enforced with an explicit seat-filling rule, plus a
budget frontier over the share of small places;
(iii) the teaching staff's subgroup-level allocation as a benchmark;
(iv) placebo and semi-synthetic tests of whether gains are real.
Preview of findings: how many small places a school has matters far more than who gets
them (+0.20 SD from all-regular to all-small vs at most about +0.03 to +0.05 SD from
targeting); the one real difference between students is gender (boys gain more), and
targeting on it nearly shuts girls out of small classes.}

% =====================================================================================
\section{Project STAR Data}  % ~0.75 page

\subsection{Study Design and Treatment Arms}
\placeholder{Random assignment of students and teachers to class types WITHIN schools.
Unit: student. We use the kindergarten cohort and the class type in the student's first
STAR year (the randomized one; intent-to-treat), because later switches were not random.
Data: AER STAR file (Rdatasets), 11,598 students. Notes sections 2--3.}

\subsection{Variables, Outcomes, and Sample Construction}
\placeholder{Outcome: mean of reading and math z-scores at the end of kindergarten,
standardized against regular and regular-with-aide students (SD units). Features known
before assignment: gender, ethnicity (white / Black / other), age at entry, free lunch,
school locale (rural / suburban / urban / inner-city). Teacher variables excluded because
they depend on the class assigned. Sample: 6,325 kindergarten entrants; dropped 539 with no
outcome (missing at the same rate, 8.4--8.6\%, in every class type) and 34 in school 14
(no regular class, so no overlap); final 5,752 students in 78 schools. A few missing
features filled with the median, with flags. Propensities: each class type's share within
the school (minimum 0.161). Table~\ref{tab:descriptives}: the arms are balanced, as
randomization implies.}
% NOTE: grade-3 outcome is 51% missing (attrition); mention as a limitation only.

\begin{table}[htbp]
\centering
\caption{Kindergarten analysis sample by assigned class type (source:
\texttt{figures/table1\_data\_summary.csv}). Outcomes in SD units relative to regular
and regular-with-aide students.}
\label{tab:descriptives}
\begin{tabular}{lcccc}
\toprule
 & Regular & Small & Regular + aide & All \\
\midrule
Students        & 2{,}005 & 1{,}725 & 2{,}022 & 5{,}752 \\
Female (\%)     & 48.7 & 48.5 & 48.7 & 48.6 \\
Free lunch (\%) & 47.4 & 46.9 & 49.4 & 47.9 \\
Black (\%)      & 31.7 & 30.6 & 32.6 & 31.7 \\
Age at entry (years) & 5.42 & 5.43 & 5.43 & 5.43 \\
Rural school (\%) & 50.0 & 46.6 & 45.5 & 47.4 \\
Outcome (SD)    & 0.00 & 0.18 & 0.01 & 0.06 \\
\bottomrule
\end{tabular}
\end{table}

% =====================================================================================
\section{Methods}  % ~1.75 pages

\subsection{Estimating Counterfactual Rewards}
\placeholder{Each student is observed in one class type only, so we estimate a reward
for every student under every class type. Direct Method (DM): arm-specific outcome models
$\widehat f_t$ (random forest; linear as a check). Doubly robust (DR), our main
estimator, adds an inverse-propensity correction. Cross-fitting: training rewards come
from outcome models that never saw the student's own fold; test rewards from models fit
only on training schools. Single-student DR rewards are very noisy (spread about 1.8 SD
vs about 0.3 for DM): they are for averaging, not individual effects.}

Let $x_i$ be student $i$'s features, $T_i\in\treatmentset=\{\mathrm{reg},\mathrm{small},
\mathrm{aide}\}$ the assigned class type, and $y_i$ the observed outcome. The doubly
robust reward is
\begin{equation}
\widehat{\Gamma}_{it}
=\widehat{f}_{t}(x_i)
+\frac{\mathbb{1}\{T_i=t\}}{\widehat{p}_{t}(x_i)}
\bigl(y_i-\widehat{f}_{t}(x_i)\bigr).
\label{eq:dr}
\end{equation}

\subsection{Optimal Policy Trees}
\placeholder{A policy tree splits on features and prescribes a class type in each leaf.
Fit with IAI's OptimalTreePolicyMaximizer \cite{amram2022,iai}, depth 1--3 tuned by
GridSearch, minimum leaf size 50. IAI uses local search, so the tree is not certified
optimal. For allocation we lower the small (and aide) reward by its mean training gain, so
every split separates students who gain more than average from those who gain less; each
leaf's mean training gain ($\widehat\Gamma_{\mathrm{small}}-\widehat\Gamma_{\mathrm{reg}}$)
becomes a priority score.}

\begin{equation}
\max_{\pi\in\Pi_D}\; \sum_{i=1}^{n}\widehat{\Gamma}_{i,\pi(x_i)}.
\label{eq:opt}
\end{equation}

\subsection{School-Level Capacity Constraints}
For school $s$ with students $S_s$ and available places $B_s^{\mathrm{small}}$,
$B_s^{\mathrm{aide}}$:
\begin{equation}
\sum_{i\in S_s}\mathbb{1}\{\pi(x_i)=\mathrm{small}\}\leq B_s^{\mathrm{small}},\qquad
\sum_{i\in S_s}\mathbb{1}\{\pi(x_i)=\mathrm{aide}\}\leq B_s^{\mathrm{aide}}.
\label{eq:capacity}
\end{equation}
\placeholder{A tree alone cannot guarantee (\ref{eq:capacity}) in new schools: used as is,
it breaks a quota in 53--71 of the 78 schools. We therefore fill seats school by school
with the same rule for every policy: small places go to the $B_s^{\mathrm{small}}$
students with the highest priority, then aide places, the rest regular. Two scenarios,
both ASSUMED, not real budgets:
(a) budget frontier: $B_s^{\mathrm{small}}=\mathrm{round}(q\,n_s)$, $q=0,0.1,\dots,1$,
no aide places;
(b) capacity scale: $B_s = \mathrm{round}(c\times$ the school's actual STAR seats$)$,
$c=0.5,\dots,1.5$; at $c=1$ every school keeps its real classes and only who sits in them
changes. This is a seat-filling heuristic, not an exact MIP with (\ref{eq:capacity})
inside the fit (future work). Notes 12c, 13d.}
% NOTE: the earlier lambda cost penalty only controls the budget on average, not per
% school. Mention it in one sentence at most, or leave it out (decide with Daniel).

\subsection{Subgroup-Level Allocation (Teaching Staff Benchmark)}
\placeholder{Following the staff's suggestion: groups = every combination of gender x
free lunch x ethnicity x locale x three age bands (144 possible; median about 15 training
students per group). Each group's priority = its mean training gain; the same seat-filling
rule as OPT. With one scarce resource per school, filling by group gain solves the staff's
small linear program exactly. Unlike OPT, the groups are fixed in advance, not learned.
Notes 14a.}

\subsection{Checking That Gains Are Real}
\placeholder{Placebo: shuffle class types within each school (20 shuffles) and rerun the
whole pipeline; any "gain" is chance. p-value of a real result $=(1+\#\text{placebos at or
above})/21$. Semi-synthetic test: add a known extra small-class benefit to one group
(labelled SYNTHETIC) and check the tree finds it. Uncertainty: school-level bootstrap
(2,000 draws, trees held fixed).}

% =====================================================================================
\section{Experiments and Results}  % ~3.5 pages

\subsection{Experimental Setup and Baselines}
\placeholder{Five folds of whole schools (seed 15): each school is a test school once and
never used to fit its own policy. Policy value on held-out students:}
\begin{equation}
\widehat V(\pi)=\frac{1}{n_{\mathrm{test}}}
\sum_{i\in D_{\mathrm{test}}}\widehat{\Gamma}^{\mathrm{test}}_{i,\pi(x_i)}.
\label{eq:policyvalue}
\end{equation}
\placeholder{Baselines. Unconstrained references: all-regular ($-0.014$), all-small
($0.200$), all-aide (about the same as all-regular). Same-budget comparisons (same seats
in every school): random allocation (exact expected value), free-lunch-first, subgroup
table. We report every policy as its gain over random at the same budget.}

\subsection{Budget Frontier: Does Targeting Beat Random?}
\placeholder{Main result (Fig.~\ref{fig:frontier}, notes 12d--e). OPT beats random by
+0.032 SD at $q=0.3$ (95\% CI [+0.004, +0.062]) and +0.032 at $q=0.4$ ([+0.001, +0.065]);
elsewhere the interval covers 0. Against the 20 placebos the best share gives p = 0.095
(DR rewards); trees trained on linear rewards reach +0.049 (p = 0.048, the smallest
possible with 20 placebos). Free-lunch-first is slightly BELOW random at every share:
free-lunch students gain the same as others. Takeaway: going from 0\% to 100\% small
classes is worth +0.21 SD; choosing who gets them is worth at most about +0.03 to +0.05,
at the edge of chance.}

\begin{figure}[htbp]
\centering
\fig{figures/fig2_frontier.png}{0.85\linewidth}{budget frontier}
\caption{Gain over random allocation (held-out DR value, SD) vs share of students given
small-class places in each school. Bands: 95\% school-bootstrap intervals; grey: placebo
5--95\% range.}
\label{fig:frontier}
\end{figure}

\subsection{Who Benefits: The Learned Trees}
\placeholder{Every frontier tree splits on gender, in all 5 folds and for all four reward
types (placebo trees: 29 of 100 folds). Held-out small-class gain: boys +0.306 SD, girls
+0.117 (gap +0.19, 95\% CI [+0.07, +0.31]); raw means (+0.21) and a regression with
school fixed effects (+0.135) agree. No other difference is significant: free lunch 0.215
vs 0.212. Fig.~\ref{fig:trees}: fold-0 tree, "boy?" with leaf gains +0.29 vs +0.12.
Without gender and ethnicity the tree splits on locale (+0.25 vs +0.11), a SCHOOL-level
feature, which cannot rank students inside a school, so within-school filling is
effectively random. The aide arm adds nothing measurable. Notes 12f, 13c.}
% NOTE: Daniel's notes say not to cite literature on gender differences without checking
% it. Present the gender gap as our finding.

\begin{figure}[htbp]
\centering
\fig{figures/fig5_trees.png}{0.85\linewidth}{example trees}
\caption{Fold-0 policy trees: with all features (left) and without gender and ethnicity
(right). Leaves show the mean estimated small-class gain.}
\label{fig:trees}
\end{figure}

\subsection{Sensitive Features and Fairness}
\placeholder{All of the targeting value comes from gender. Without gender and ethnicity,
OPT is no better than random (best share +0.005, p = 0.86). With them, at $q=0.3$ the
policy gives small classes to 53\% of boys but 5\% of girls (random: 30\% each); boys gain
+0.107 [+0.063, +0.154] over random, girls LOSE $-0.047$ [$-0.085$, $-0.007$]
(Fig.~\ref{fig:gender}). Without the sensitive features every group gets about 30\% and no
group gains or loses. Notes 13c.}

\begin{figure}[htbp]
\centering
\fig{figures/fig4_gender_shares.png}{0.8\linewidth}{gender shares}
\caption{At 30\% small-class places: share of boys and girls given small classes, and
gain over random within each gender, for OPT with and without gender/ethnicity.}
\label{fig:gender}
\end{figure}

\subsection{Per-School Capacities}
\placeholder{(Fig.~\ref{fig:capacity} or a small table; notes 13e.) OPT minus random at
$c=0.5,0.75,1,1.25,1.5$: +0.028, +0.022, +0.006, +0.014, +0.004; inside the placebo's 95th
percentile at every scale. At the real STAR seat counts ($c=1$), rearranging who sits in
which class gains +0.006 SD. Every feasible policy uses exactly its quota in every school;
the tree's own prescriptions would not (53--71 schools over quota).}

\begin{figure}[htbp]
\centering
\fig{figures/fig3_capacity.png}{0.8\linewidth}{capacity scenarios}
\caption{Gain over random allocation under per-school capacities, scaled from each
school's observed STAR seats ($c=1$: actual seats).}
\label{fig:capacity}
\end{figure}

\subsection{OPT vs Subgroup-Level Allocation}
\placeholder{Answer to the staff's question (notes 14b). With DR rewards, the 144-group
table overfits (about 15 noisy rewards per group) and lands at random; OPT beats it by
+0.034 [+0.002, +0.067] at $q=0.3$: the shallow tree gives up nothing for
interpretability and generalizes better. With smoother DM-linear rewards the table and the
tree tie (both about +0.05 over random, p = 0.048), recovering the same gender gap.
Without gender and ethnicity, and under per-school capacities, nothing beats random.}

\subsection{Robustness}
\placeholder{One paragraph. (i) Reward estimators: DM vs DR, random forest vs linear, all
scored with the same DR test rewards: same picture. (ii) Semi-synthetic (SYNTHETIC): a
planted extra +0.3 SD for free-lunch students is found in 5 of 5 folds (leaf gains 0.20 vs
0.51; OPT ties the ideal rule); +0.2 in 3 of 5; +0.1 in 1 of 5. So the method can find
differences of about 0.2--0.3 SD, and the real data show none besides gender. (iii) The
earlier cost-penalty sweep and its placebos (OPT minus random $-0.032$ to $+0.015$) agree.
Notes 11, 12b.}
% NOTE: No training-set (in-sample) values were computed, so the old outline's
% "compare training and held-out performance" is dropped.

% =====================================================================================
\section{Discussion}  % ~0.9 page

\subsection{Implications for Resource Allocation}
\placeholder{The number of small places matters far more than who gets them. Targeting
on observable traits adds little; the one signal (gender) comes with a fairness cost we
would not recommend accepting, so our recommended policy is the one without sensitive
features, which in practice is close to random (lottery-like) allocation within schools.
Aides add nothing measurable. A shallow tree does at least as well as the fine subgroup
table, so interpretability costs nothing here. Seat limits need an explicit filling rule;
a tree alone ignores them.}

\subsection{Limitations}
\placeholder{Kindergarten outcome only (grade 3 is 51\% missing); few, coarse features
(no baseline test score); noisy DR rewards and borderline p-values with 20 placebos;
bootstrap holds trees fixed; capacities are assumed scenarios and counted in seats, not
whole classrooms or staff; seat filling is a heuristic, not an exact MIP; no locked test
set beyond 5-fold cross-validation; results may not generalize beyond 1980s Tennessee.}
% NOTE: if the exact MIP (task 19) gets done, move it from Limitations to Methods/Results.

% =====================================================================================
\section{Conclusion}  % ~0.25 page
\placeholder{Answer the research question in one paragraph: in STAR kindergarten,
observable student traits give little to target on; the only one that matters (gender)
would make allocation unfair; capacity, not targeting, drives achievement. Next step:
richer pre-assignment data (e.g. baseline skills) and an exact MIP with seat limits inside
the tree fit.}

% =====================================================================================
\begin{thebibliography}{9}
\bibitem{amram2022} M. Amram, J. Dunn, and Y. D. Zhuo. ``Optimal Policy Trees.''
\emph{Machine Learning}, 111(7):2741--2768, 2022.
\bibitem{iai} Interpretable AI, LLC. \emph{Interpretable AI Documentation}.
\url{https://www.interpretable.ai}.
\bibitem{krueger1999} A. B. Krueger. ``Experimental Estimates of Education Production
Functions.'' \emph{The Quarterly Journal of Economics}, 114(2):497--532, 1999.
\bibitem{nye2000} B. Nye, L. V. Hedges, and S. Konstantopoulos. ``The Effects of Small
Classes on Academic Achievement: The Results of the Tennessee Class Size Experiment.''
\emph{American Educational Research Journal}, 37(1):123--151, 2000.
\bibitem{whitehurst2011} G. J. Whitehurst and M. M. Chingos. \emph{Class Size: What
Research Says and What It Means for State Policy}. Brookings Institution, 2011.
\end{thebibliography}

\section*{Contributions}
\placeholder{Agree the wording together. For reference, from the repo history: Daniel
built the data download and cleaning, the reward and tree pipeline, the placebo, budget
frontier, capacity, sensitive-feature and subgroup analyses, and the figures; Adi
reproduced the pipeline, ran the semi-synthetic test and the DM/linear comparisons, fixed
the linear-model reward bug, and wrote the report.}

\end{document}
